\documentclass[11pt]{article}

\usepackage[a4paper,margin=2.4cm]{geometry}
\usepackage{amsmath,amssymb}
\usepackage{booktabs}
\usepackage{siunitx}
\usepackage[hidelinks]{hyperref}
\usepackage{enumitem}

\newcommand{\Rd}{R_{\mathrm d}}
\newcommand{\Rp}{R_{\mathrm p}}

\title{Spread-Out Bragg Peak (SOBP) field design\\
for the Stage-A proton-transport simulation}
\author{Method note for review}
\date{\today}

\begin{document}
\maketitle

\begin{abstract}
This note documents the method used to build a Spread-Out Bragg Peak (SOBP) in
the Stage-A simulation, before it is implemented. A single mono-energetic pencil
deposits a sharp Bragg peak and paints a thin track, giving an unrealistically
small integral activity. A clinical field instead spreads dose uniformly over a
target volume, which is what fixes the realistic absolute $\beta^{+}$ yield (and
lets us cross-check against Parodi \emph{et al.}\ 2008, Table~2). We construct
the depth component of such a field analytically: a weighted superposition of
pristine Bragg peaks whose weights flatten the dose over the target's depth
extent. The energies and weights are computed in Python
(\texttt{field\_design/sobp.py}) and written to a layer table that the Geant4
gun samples from; the resulting plateau is then verified directly against the
simulated depth--dose curve. Lateral spreading and dose normalisation are
separate steps, noted at the end.
\end{abstract}

\section{Goal and scope}

We want a proton field that deposits a \textbf{flat dose plateau} across the
target's depth extent $[\Rp,\Rd]$ (proximal to distal range), so that a
prescribed dose can be painted \emph{uniformly} over the target. This note
covers only the \textbf{depth} component (the SOBP proper). The target geometry,
its size and the prescribed dose are all \textbf{parameters}
(Sec.~\ref{sec:params}); nothing here is hard-wired to the Parodi reference
numbers.

\section{Pristine Bragg peak and the range--energy relation}
\label{sec:pristine}

A mono-energetic proton of energy $E$ stops at a well-defined depth, its
\textbf{range} $R$. To good approximation the two are related by the
Bortfeld power law~\cite{bortfeld1996},
\begin{equation}
R(E) = \alpha\,E^{\,p},
\label{eq:range}
\end{equation}
with, in water, $\alpha=\SI{2.2e-3}{cm\,MeV^{-p}}$ and $p=1.77$. For another
homogeneous material we scale the range by the (mass) stopping-power ratio,
which to first order is the inverse mass-density ratio,
$R_{\mathrm{mat}}\simeq R_{\mathrm{water}}/\rho_{\mathrm{rel}}$
($\rho_{\mathrm{rel}}\approx1.04$ for brain). The energy needed to reach a depth
$R$ is the inverse of Eq.~\eqref{eq:range}, $E(R)=(R/\alpha)^{1/p}$.

Equation~\eqref{eq:range} is used only to choose the \emph{layer energies}; the
actual peak shapes come from Geant4, and the plateau is verified there
(Sec.~\ref{sec:verify}), so small inaccuracies in $\alpha,p$ are corrected by
the verification step.

A useful consequence of Eq.~\eqref{eq:range}: a proton that has travelled to
depth $z<R$ has residual range $R-z$ and residual energy
$\propto (R-z)^{1/p}$, so its mass stopping power (dose per unit length) scales as
\begin{equation}
S(z;R)\ \propto\ \frac{\mathrm dE}{\mathrm dz}\ \propto\ (R-z)^{\,1/p-1}.
\label{eq:stop}
\end{equation}
This is the only feature of the peak shape the weight derivation needs.

\section{SOBP as a weighted superposition}

We place $N{+}1$ pristine peaks at ranges $R_i$ spanning the target depth,
\begin{equation}
R_i = \Rp + \frac{i}{N}\,(\Rd-\Rp),\qquad i=0,\dots,N,
\end{equation}
($R_0=\Rp$ proximal, $R_N=\Rd$ distal), with relative fluence weights $w_i\ge0$.
The total depth--dose is their superposition,
\begin{equation}
D_{\mathrm{SOBP}}(z)=\sum_{i\,:\,R_i\ge z} w_i\,S(z;R_i).
\label{eq:super}
\end{equation}
The gun realises Eq.~\eqref{eq:super} stochastically: each primary proton is
assigned to layer $i$ with probability $\propto w_i$ and given energy
$E_i=E(R_i)$.

\section{Analytic flattening weights}
\label{sec:weights}

We want $D_{\mathrm{SOBP}}(z)=\text{const}$ for $z\in[\Rp,\Rd]$. Treating the
range as continuous with fluence density $\Phi(R)$, Eq.~\eqref{eq:super} with
\eqref{eq:stop} becomes the Abel-type integral
\begin{equation}
D(z)=\int_z^{\Rd}\Phi(R)\,(R-z)^{1/p-1}\,\mathrm dR \stackrel{!}{=}\text{const}.
\end{equation}
Its solution---the classic result of Bortfeld \& Schlegel~\cite{bortfeld1996}---is
\begin{equation}
\boxed{\;\Phi(R)\ \propto\ (\Rd-R)^{\,1/p-1}\;}\qquad R\in[\Rp,\Rd],
\label{eq:phi}
\end{equation}
plus a finite reinforcement at the proximal edge $R=\Rp$ that fills the proximal
corner of the plateau. Since $1/p-1\approx-0.435<0$, the weight rises towards the
\emph{distal} peak (which no deeper peak reinforces) and diverges integrably at
$R=\Rd$. Integrating Eq.~\eqref{eq:phi} over each range bin
$[R_i-\tfrac{\Delta}{2},\,R_i+\tfrac{\Delta}{2}]$ (with $\Delta=(\Rd-\Rp)/N$)
gives the finite discrete weights actually used,
\begin{equation}
w_i\ \propto\ (\Rd-R_i+\tfrac{\Delta}{2})^{1/p}-(\Rd-R_i-\tfrac{\Delta}{2})^{1/p},
\label{eq:wi}
\end{equation}
with the distal bin truncated at $\Rd$ and an added proximal term for $i=0$.
The weights are normalised to sum to one and stored as a probability table.

\paragraph{Why analytic (vs.\ fit-to-simulation).} Equation~\eqref{eq:wi} needs
no simulation and is the standard, well-validated construction. We adopt it as
the baseline; if the verified plateau (Sec.~\ref{sec:verify}) is not flat enough,
the fallback is to replace the analytic $S(z;R_i)$ with Geant4-simulated pristine
peaks and solve Eq.~\eqref{eq:super} for $w_i$ by non-negative least squares.

\section{Parameters}
\label{sec:params}

Everything that defines the target and the field is a parameter:
\begin{center}
\begin{tabular}{lll}
\toprule
symbol & meaning & Parodi-standard value \\
\midrule
$d_{\mathrm{prox}},d_{\mathrm{dist}}$ & target proximal/distal depth & $\approx 5.5,\ \SI{10.5}{cm}$ \\
$\varnothing_{\mathrm t}$ & target (tumour) diameter & \SI{6}{cm} \\
$L_{\mathrm t}=d_{\mathrm{dist}}-d_{\mathrm{prox}}$ & target length (modulation width) & \SI{5}{cm} \\
$D$ & prescribed dose to the target & \SI{1}{Gy} \\
$N{+}1$ & number of energy layers & $\sim 15$--$25$ \\
material & phantom material & \texttt{G4\_BRAIN\_ICRP} \\
\bottomrule
\end{tabular}
\end{center}
The depths map to ranges $\Rp,\Rd$ in the chosen material via Sec.~\ref{sec:pristine}.

\section{Algorithm and output}

\begin{enumerate}[nosep]
\item From $(d_{\mathrm{prox}},d_{\mathrm{dist}},N,\text{material})$ compute
$\Rp,\Rd$ and the layer ranges $R_i$.
\item Energies $E_i=E(R_i)$ from Eq.~\eqref{eq:range}; weights $w_i$ from
Eq.~\eqref{eq:wi}; normalise $\sum_i w_i=1$.
\item Write \texttt{data/sobp\_layers.csv} with columns
\texttt{energy\_MeV,\,weight}, and a sanity plot of the analytic SOBP.
\item The Geant4 gun loads the table and samples one layer per primary.
\end{enumerate}

\section{Verification}
\label{sec:verify}

We run Stage~A with the SOBP gun and inspect the simulated depth--dose
(\texttt{depth\_dose.csv} via \texttt{bragg\_profile.py}). Acceptance: the dose
is flat to within a few percent across $[\Rp,\Rd]$ with a sharp distal fall-off
at $\Rd$. If not, increase $N$ or fall back to the least-squares weights above.

\section{Out of scope here (next steps)}

\begin{itemize}[nosep]
\item \textbf{Lateral spreading:} uniform fluence over the target cross-section
(diameter $\varnothing_{\mathrm t}$), turning the depth plateau into a dose box.
\item \textbf{Dose normalisation:} score dose in the target volume and scale the
number of primaries so the target dose equals $D$; this fixes the absolute $P_j$.
\item \textbf{Cross-check:} compare the integral $\beta^{+}$ yields to Parodi
Table~2 ($\sim\!\SI{1.8e8}{}$ decays per Gy, head).
\end{itemize}

\begin{thebibliography}{9}
\bibitem{bortfeld1996} T.~Bortfeld and W.~Schlegel,
\emph{An analytical approximation of depth--dose distributions for therapeutic
proton beams}, Phys.\ Med.\ Biol.\ \textbf{41}, 1331 (1996); see also
T.~Bortfeld, Med.\ Phys.\ \textbf{24}, 2024 (1997).
\end{thebibliography}

\end{document}
